\subsection{Formalizing Factual Knowledge}
\label{subsec:definitions_maintext}

\paragraph{Fact sets and storage.}
Given key embeddings $\mathbf{K}\in \mathbb{R}^{|K| \times d}$ and value embeddings $\mathbf{V}\in\mathbb{R}^{|V| \times d}$, a \textit{fact set} is a map $f: [|\mathbf{K}|]\to [|\mathbf{V}|]$. We write $\mathbf{k}_{i}$ and $\mathbf{v}_{i}$ for the $i$th key and value embedding, respectively.

\begin{definition}[Fact storage]
\label{def:fact-storage}
A model $\mathbf{g}_\theta: \mathbb{R}^d\to \mathbb{R}^d$ \textit{stores a fact set} $f: [|\mathbf{K}|]\to [|\mathbf{V}|]$ given embeddings $\mathbf{K}$ and $\mathbf{V}$ if, for all $i \in [|\mathbf{K}|]$ and all $j \neq f(i) \in [|\mathbf{V}|]$,
\begin{equation}
    \label{eq:decoding-criterion-maintext}
    \langle \mathbf{g}_\theta(\mathbf{k}_{i}), \mathbf{v}_{f(i)} \rangle > \langle \mathbf{g}_\theta(\mathbf{k}_{i}), \mathbf{v}_{j} \rangle.
\end{equation}
\end{definition}

Notably, this definition is equivalent to correct softmax decoding in language modeling.

\begin{definition}[Margin]
\label{def:margin}
The margin of $\mathbf{g}_\theta$ on fact $i$ against competitor $j \neq f(i)$ is
\begin{equation}
    \label{eq:margin-maintext}
    \gamma_{i,j} := \langle \mathbf{g}_\theta(\mathbf{k}_{i}), \mathbf{v}_{f(i)} \rangle - \langle \mathbf{g}_\theta(\mathbf{k}_{i}), \mathbf{v}_{j} \rangle,
\end{equation}
and the minimum margin is $\gamma_{\min} := \min_{i, j \neq f(i)} \gamma_{i,j}$.
Note that storing a fact set (in the sense of~\Cref{def:fact-storage}) is equivalent to $\gamma_{\min} > 0$.
\end{definition}


\paragraph{Fact-storage cost and capacity.}
To measure parameter efficiency, we define the smallest parameter budget needed for a model class to store \emph{every} fact set on fixed embeddings.

\begin{definition}[Fact-storage cost and capacity]
\label{def:fact-storage-cost}
The \emph{fact-storage cost} of a model class $\mathbf{g}$ on embeddings $\mathbf{K}$ and $\mathbf{V}$ is the minimum parameter count needed to represent \emph{all} possible fact sets:
\begin{equation}
    W(\mathbf{g}; \mathbf{K}, \mathbf{V}) =
    \min \left\{\#(\theta)\Bigg|\;
    \begin{aligned}
        &\forall f : [|\mathbf{K}|] \to [|\mathbf{V}|], \\
        &\exists\, \theta \; \text{s.t.} \; \mathbf{g}_\theta \text{ stores } f
    \end{aligned}
    \right\}.
    \label{def:complexity}
\end{equation}
The corresponding \emph{fact-storage capacity} is the maximum number of facts storable with a fixed parameter budget.
\end{definition}

\begin{theorem}[Information-theoretic lower bound]
\label{thm: info_bounds_capacity-const}
Assuming a constant number of bits per parameter, the fact-storage cost of embeddings $\mathbf{K}$ and $\mathbf{V}$ for \emph{any} model class $\mathbf{g}$ satisfies
\[
W(\mathbf{g}; \mathbf{K}, \mathbf{V}) = \Omega(|\mathbf{K}|\log [|\mathbf{V}|]).
\]
\end{theorem}
See Appendix~\ref{sec:info_theory_bound} for proof.


\subsection{Model Classes}
\label{subsec:model_classes}
In this work, we study two model classes: gated one-hidden-layer MLPs and Hebbian memories~\citep{10.1109/TC.1972.5008975, doi:10.1073/pnas.79.8.2554, bubeck2020networksizeweightssize, cabannes2024scalinglawsassociativememories, nichani2024understandingfactualrecalltransformers}.

\textbf{MLPs.} We consider models $\mathbf{g}_\theta:\mathbb{R}^d\to\mathbb{R}^{d_v}$ of the form
\begin{equation}
\label{eq:gated-mlp-class}
\text{MLP}(\mathbf{x}) =
\mathbf{g}_\theta(\mathbf{x})=\mathbf{B}\!\left((\mathbf{A}\mathbf{x})\odot \sigma(\mathbf{G}\mathbf{x})\right),
\end{equation}
where $\mathbf{x}\in\mathbb{R}^d$, $\mathbf{A},\mathbf{G}\in\mathbb{R}^{m\times d}$, and $\mathbf{B}\in\mathbb{R}^{d_{v}\times m}$.
This family includes SwiGLU-style MLPs~\citep{shazeer2020gluvariantsimprovetransformer} used in modern language models~\citep{yang2025qwen3technicalreport, deepseekai2025deepseekv3technicalreport, dubey2024llama}. Our explicit construction (\Cref{sec:construction}) uses $\sigma=\mathrm{id}$.

\textbf{Hebbian memories.} These linear models are maps $\mathbf{g}_{\mathbf{W}}:\mathbb{R}^d\to\mathbb{R}^{d_v}$ of the form $\mathbf{g}_{\mathbf{W}}(\mathbf{x})=\mathbf{W}\mathbf{x}$, where $\mathbf{x}\in\mathbb{R}^d$ and $\mathbf{W}\in\mathbb{R}^{d_{v}\times d}$. A fact set is stored by taking
\begin{equation}
\label{eq:hebbian-memory-class}
\mathbf{W}=\sum_{j=1}^{|\mathbf{K}|}\mathbf{v}_{f(j)}\mathbf{k}_{j}^\top.
\end{equation}
Section~\ref{sec:3} shows that MLPs can be recast as Hebbian memories in a kernel feature space.


\subsection{Related Work}
\label{sec:related_work}

The two closest works to ours are \citet{nichani2024understandingfactualrecalltransformers} and \citet{zhong2025understandingtransformerperspectiveassociative}. \citet{nichani2024understandingfactualrecalltransformers} gave the first explicit construction of fact-storing MLPs and showed near-optimal fact-storage capacity, up to a polylogarithmic factor, but their analysis is restricted to isotropic embeddings and studies only the separability condition $\gamma_{\min}>0$. \citet{zhong2025understandingtransformerperspectiveassociative} developed a unified associative-memory view of attention and MLPs
, but focused on average retrieval fidelity rather than worst-case margins. In contrast, we derive explicit margin bounds beyond isotropy, identify margin bounded away from zero as the condition for Transformer usability, and show that constructed MLPs can be integrated into Transformer blocks for factual recall.

Additional discussion of probing, editing, and scaling studies of factual knowledge in language models appears in Appendix~\ref{app:related_work}.
